% Generated draft: requires mathematical and source review.
\documentclass[11pt]{article}
\usepackage[margin=25mm]{geometry}
\usepackage{fontspec}
\setmainfont{Latin Modern Roman}
\usepackage{amsmath,amssymb,mathtools,unicode-math}
\setmathfont{Latin Modern Math}
\usepackage{enumitem,needspace,xurl,hyperref}
\setlength{\parindent}{0pt}
\setlength{\parskip}{5pt}
\setlength{\emergencystretch}{3em}
\newcommand{\N}{\mathbb{N}}
\newcommand{\Z}{\mathbb{Z}}
\newcommand{\Q}{\mathbb{Q}}
\newcommand{\R}{\mathbb{R}}
\newcommand{\C}{\mathbb{C}}
\newcommand{\F}{\mathbb{F}}
\newcommand{\A}{\mathbb{A}}
\newcommand{\PP}{\mathbb{P}}
\newcommand{\E}{\mathbb{E}}
\newcommand{\eps}{\varepsilon}
\newcommand{\dd}{\,\mathrm{d}}
\begin{document}
\section*{UnsolvedMath ID: 30001075}
\textbf{Source problem code:} catalogue entry
\subsection{Definitions and mathematical statement}
Let \(A,B,C\subseteq\R^3\) be pairwise disjoint convex sets:
\[
 A\cap B=A\cap C=B\cap C=\varnothing.
\]
For the usual closed convex-body interpretation, an affine line \(\ell\) is tangent to \(A\) when \(\ell\cap A\ne\varnothing\) and \(\ell\cap\operatorname{int}(A)=\varnothing\). This convention allows contact along a boundary segment or face, not just a unique smooth contact point. The extraction does not require closedness, compactness, or nonempty interior, so tangency for general convex sets requires an additional source convention.

With the intended tangency relation, put
\[
 \mathcal T(A,B,C)=\{\ell:\ell\text{ is an affine line tangent to }A,B,C\},
 \qquad
 S(A,B,C)=\bigcup_{\ell\in\mathcal T(A,B,C)}\ell.
\]
Equivalently, \(x\in S(A,B,C)\) if some \(v\ne0\) makes the line \(x+\R v\) tangent to all three sets. The point \(x\) can lie anywhere on the common tangent line, not only at a tangency point.

Conjecture 4 asks whether
\[
 \forall A,B,C\text{ as above}\quad
 \exists Z\subseteq\R^3\text{ Lebesgue measurable}:\quad
 S(A,B,C)\subseteq Z,\qquad\mathcal L^3(Z)=0.
\]
Equivalently, the outer three-dimensional Lebesgue measure of \(S(A,B,C)\) is zero. This states the source's null-set containment without assuming a priori measurability of the tangent locus. It is not a claim that the line family is countable or that \(S\) is a smooth surface; uncountable unions of lines need not be null merely because each line is null.

The source describes this as a weaker visual-event-surface conjecture. No stronger dimension, regularity, or surface decomposition statement is supplied. The assessment dated 2026-08-21 says the extraction contains only the label of that conjecture, whereas the supplied original and clean text explicitly define the locus and assert null-set containment. This mismatch is retained as a source-review issue; the assessment supplies no resolution.
\subsection{Short English statement}
Must the union of all lines tangent to three pairwise disjoint convex sets have zero outer volume in three-dimensional space? The conjecture concerns every point on those lines. The exact convex-set and tangency conventions remain to be checked.
\subsection{Sources}
\begin{itemize}
\item Local UnsolvedMath registry, permanent ID \(30001075\), exact source code \texttt{catalogue entry}, proposed year 2008; no individual proposer is supplied. The original calls this Conjecture 4.
\item \emph{Discrete Geometry (2009)}, Oberwolfach Reports, DOI 10.4171/owr/2008/44, pp. 2477--2556. \url{https://doi.org/10.4171/owr/2008/44}.
\item The supplied \texttt{corrected\_verified} note calls the common-tangent-locus conjecture complete. The open label and assessment checked 2026-08-21 are supplied dated curation, not independently verified; its label-only description conflicts with the explicit supplied statement and requires review.
\end{itemize}
\end{document}
